\documentclass[11pt, a4paper, oneside]{article}
\usepackage[ngerman]{babel}
\usepackage{worksheet}

\begin{document}
    \author{L. Bung}
    \title{Lokale Änderungsrate und Ableitungsfunktionen}
    \subject{Mathematik}
    \class{FHR2}

    \maketitle

    \begin{hint}{Von der mittleren zur lokalen Änderungsrate}
        \setstretch{1}
        \setlength{\parskip}{.6em}
        \begin{minipage}[c]{6.3cm}
            \centering
            \begin{tikzpicture}
                \begin{axis}[axis lines=center, x=1cm, y=1cm, grid=major,
                    xlabel={$x$}, ylabel={$y$}, xlabel style={above right}, ylabel style={above right},
                    xmin=-0.5, xmax=3, ymin=-1.5, ymax=4.5, xtick={-1,...,3}, ytick={-1,...,4}, clip=false]
                    \addplot[domain=-0.5:2.12, red, thick, samples=100]{x^2};
                    \addplot[domain=0.2:2.1, blue, thick]{3*x-2};
                    \addplot[domain=0.1:2.6, ForestGreen, thick, dashed]{2*x-1};
                    \draw[dotted, thick] (1,1) -- (2,1) -- (2,4);
                    \node[below] at (1.5,1) {\small $h$};
                    \node[right, fill=white, inner sep=1pt] at (2.05,2.2) {\scriptsize $f(x_0+h)-f(x_0)$};
                    \fill (1,1) circle(2pt) node[above left]{$P$};
                    \fill (2,4) circle(2pt) node[right]{$Q$};
                    \node[red, fill=white, inner sep=1pt] at (0.85,3.4) {$f(x)=x^2$};
                    \node[blue, fill=white, inner sep=1pt] at (2.75,1.45) {\small Sekante};
                    \node[ForestGreen, fill=white, inner sep=1pt] at (2.75,0.95) {\small Tangente};
                \end{axis}
            \end{tikzpicture}
        \end{minipage}
        \hfill
        \begin{minipage}[c]{8.7cm}
            \textbf{Differenzenquotient (mittlere Änderungsrate)}\\
            Die Steigung der \textcolor{blue}{Sekante} durch die Punkte $P(x_0 \mid f(x_0))$ und $Q(x_0+h \mid f(x_0+h))$ heißt \textit{Differenzenquotient}:
            \[
                m_s = \frac{f(x_0+h)-f(x_0)}{h}
            \]
            Sie gibt an, wie stark sich $f$ \textit{im Mittel} zwischen $x_0$ und $x_0+h$ ändert.
        \end{minipage}

        \textbf{Beispiel:} Für $f(x)=x^2$ gilt an einer beliebigen Stelle $x_0$:
        \[
            m_s = \frac{(x_0+h)^2-x_0^2}{h}
                = \frac{x_0^2+2x_0h+h^2-x_0^2}{h}
                = \frac{2x_0h+h^2}{h}
                = \frac{h\cdot(2x_0+h)}{h}
                = 2x_0+h
        \]
        An der Stelle $x_0=1$ ergibt sich also $m_s = 2+h$:

        \begin{center}
            \begin{tabular}{|c|c|c|c|c|c|}
                \hline
                $h$ & $1$ & $0{,}5$ & $0{,}1$ & $0{,}01$ & $0{,}001$\\
                \hline
                $m_s = 2+h$ & $3$ & $2{,}5$ & $2{,}1$ & $2{,}01$ & $2{,}001$\\
                \hline
            \end{tabular}
        \end{center}

        \textbf{Limes:} Je kleiner $h$ wird, desto näher rückt $Q$ an $P$ heran und desto mehr nähert sich die Sekante der \textcolor{ForestGreen}{Tangente} in $P$ an.
        Den Wert, dem sich ein Term für immer kleinere $h$ annähert, schreibt man mit dem \textit{Limes} (lat.\ \glqq Grenze\grqq):
        \[
            \lim_{h \to 0} (2x_0+h) = 2x_0 \qquad \text{(gelesen: \glqq Limes für $h$ gegen $0$\grqq)}
        \]

        \textbf{Differenzialquotient (lokale Änderungsrate):} Der Grenzwert des Differenzenquotienten heißt \textit{Differenzialquotient}.
        Er gibt die Steigung der Tangente an der Stelle $x_0$ an:
        \[
            f'(x_0) = \lim_{h \to 0} \frac{f(x_0+h)-f(x_0)}{h}
        \]
        Für $f(x)=x^2$ gilt also $f'(x_0)=2x_0$, z.\,B.\ $f'(1)=2$, $f'(3)=6$ oder $f'(-2)=-4$.
        Da dies für \textit{jede} Stelle gilt, erhält man eine neue Funktion: $f'(x)=2x$.
    \end{hint}


    \task{Änderungsrate in bestimmten Punkten}

    Gegeben ist der Graph der Funktion $f$.

    \begin{plot}[xmin=-1.5, xmax=3.5, ymin=-2.5, ymax=3.5, x=1.25cm, y=1.25cm, xtick={-1,...,3}, ytick={-2,...,3}]
        \addplot[domain=-1.4:3.55, red, thick, samples=100]{1/3*x^3-x^2+1};
        \node[red] at (1,1) {$f(x)$};
    \end{plot}

    a) Zeichnen Sie an den angegebenen Stellen die Tangente an den Graphen von $f$ ein.
    Bestimmen Sie mithilfe eines Steigungsdreiecks die lokale Änderungsrate $f'(x)$ und tragen Sie sie in die Tabelle ein.

    \begin{center}
        \begin{tabularx}{.8\textwidth}{|c|X|X|X|X|}
            \hline
            Stelle $x$ & \centering $-1$ & \centering $0$ & \centering $1$ & \centering\arraybackslash $3$\\
            \hline
            lokale Änderungsrate $f'(x)$ & & & & \\
            \hline
        \end{tabularx}
    \end{center}

    b) Nun ist umgekehrt die lokale Änderungsrate gegeben.
    Bestimmen Sie \textit{alle} Stellen $x$, an denen der Graph von $f$ diese Steigung hat.
    Falls es keine solche Stelle gibt, begründen Sie, warum.

    \begin{center}
        \begin{tabularx}{.8\textwidth}{|c|X|X|X|}
            \hline
            lokale Änderungsrate $f'(x)$ & \centering $0$ & \centering $1$ & \centering\arraybackslash $-2$\\
            \hline
            Stelle(n) $x$ & & & \\
            \hline
        \end{tabularx}
    \end{center}

    \checkered[2.5cm]


    \begin{warning}{Die Ableitungsfunktion}
        Die \textbf{lokale Änderungsrate} einer Funktion $f$ an der Stelle $x_0$ ist die Steigung der Tangente an den Graphen von $f$ im Punkt $P(x_0 \mid f(x_0))$.
        Man nennt sie auch die \textbf{Ableitung} von $f$ an der Stelle $x_0$ und schreibt $f'(x_0)$ (gelesen: \glqq $f$ Strich von $x_0$\grqq).

        Ordnet man \textit{jeder} Stelle $x$ die lokale Änderungsrate $f'(x)$ zu, erhält man eine neue Funktion: die \textbf{Ableitungsfunktion} $f'$.
        \[
            f'(x) = \lim_{h \to 0} \frac{f(x+h)-f(x)}{h}
        \]
        Der Graph von $f'$ zeigt also an jeder Stelle, wie steil der Graph von $f$ dort ist:
        \begin{itemize}
            \item $f'(x) > 0$: Der Graph von $f$ \textbf{steigt} an der Stelle $x$.
            \item $f'(x) < 0$: Der Graph von $f$ \textbf{fällt} an der Stelle $x$.
            \item $f'(x) = 0$: Der Graph von $f$ hat an der Stelle $x$ eine \textbf{waagerechte Tangente}.
        \end{itemize}
    \end{warning}


    \task{Die Ableitungsfunktion}

    Gegeben ist der Graph der Funktion $f$.

    \begin{plot}[xmin=-2.5, xmax=6.5, ymin=-3.5, ymax=3, xtick={-2,...,6}, ytick={-3,...,2}]
        \addplot[domain=-2.5:6.5, red, thick, samples=100]{-1/4*x^2+x+1};
        \node[red] at (0.5,2.3) {$f(x)$};
    \end{plot}

    a) Bestimmen Sie mithilfe von Tangenten die lokale Änderungsrate von $f$ an den Stellen $x=-2$, $x=0$, $x=2$, $x=4$ und $x=6$.
    Zeichnen Sie anschließend für jede dieser Stellen den Punkt $(x \mid f'(x))$ in das Koordinatensystem ein.

    b) Verbinden Sie die Punkte aus a) zu einem Funktionsgraphen.
    Dies ist der Graph der Ableitungsfunktion $f'$.
    Beschreiben Sie den Verlauf des Graphen von $f'(x)$.

    \checkered[1cm]

    c) Sagen Sie mithilfe des Graphen von $f'$ voraus, welche Steigung der Graph von $f$ an der Stelle $x=3$ hat.
    Überprüfen Sie Ihre Vorhersage, indem Sie die Tangente an den Graphen von $f$ an der Stelle $x=3$ einzeichnen und ihre Steigung bestimmen.

    \checkered[3cm]


    \task{Von der Ableitungsfunktion zum Ursprungsgraphen}

    Gegeben sind jeweils links die Graphen von Ableitungsfunktionen $f'$.
    Skizzieren Sie im rechten Koordinatensystem jeweils einen möglichen Graphen der zugehörigen Funktion $f$.

    \textit{Tipp:} Überlegen Sie zuerst, wo der Graph von $f$ steigt, fällt oder eine waagerechte Tangente hat.

    \pgfplotsset{
        ursprung/.style={
            axis lines=center, x=1cm, y=1cm, grid=major,
            xmin=-3, xmax=3, ymin=-3, ymax=3,
            xtick={-2,...,2}, ytick={-2,...,2},
            tick label style={font=\footnotesize},
            xlabel={$x$}, ylabel={$y$},
            xlabel style={above right}, ylabel style={above right},
            title style={yshift=2ex},
        }
    }

    \begin{center}
        \begin{tikzpicture}
            \begin{axis}[ursprung, title={a) Graph von $f'$}]
                \addplot[domain=-3:3, blue, thick]{x-1};
            \end{axis}
        \end{tikzpicture}
        \hspace{1.5cm}
        \begin{tikzpicture}
            \begin{axis}[ursprung, title={möglicher Graph von $f$}]
                \addplot[draw=none, domain=-3:3]{0};
            \end{axis}
        \end{tikzpicture}
    \end{center}

    \begin{center}
        \begin{tikzpicture}
            \begin{axis}[ursprung, title={b) Graph von $f'$}]
                \addplot[domain=-3:3, blue, thick, samples=100]{x^3};
            \end{axis}
        \end{tikzpicture}
        \hspace{1.5cm}
        \begin{tikzpicture}
            \begin{axis}[ursprung, title={möglicher Graph von $f$}]
                \addplot[draw=none, domain=-3:3]{0};
            \end{axis}
        \end{tikzpicture}
    \end{center}

    \begin{center}
        \begin{tikzpicture}
            \begin{axis}[ursprung, title={c) Graph von $f'$}]
                \addplot[domain=-3:3, blue, thick, samples=100]{1/4*(x^3-3*x+2)};
            \end{axis}
        \end{tikzpicture}
        \hspace{1.5cm}
        \begin{tikzpicture}
            \begin{axis}[ursprung, title={möglicher Graph von $f$}]
                \addplot[draw=none, domain=-3:3]{0};
            \end{axis}
        \end{tikzpicture}
    \end{center}

    d) Vergleichen Sie Ihre Graphen mit denen Ihres Sitznachbarn bzw.\ Ihrer Sitznachbarin.
    Sind Ihre Graphen gleich?
    Erklären Sie, warum es zu einer Ableitungsfunktion mehrere mögliche Funktionen $f$ geben kann.

    \checkered[4cm]


    \task{Zuordnungen}

    Oben sind die Graphen von vier Funktionen A bis D abgebildet, unten die Graphen von vier Ableitungsfunktionen 1 bis 4.
    Ordnen Sie jeder Funktion die passende Ableitungsfunktion zu und begründen Sie Ihre Zuordnung.

    \pgfplotsset{
        zuordnung/.style={
            xmin=-2.5, xmax=2.5, ymin=-2.5, ymax=2.5,
            x=.6cm, y=.6cm,
            xtick={-2,...,2}, ytick={-2,...,2},
            tick label style={font=\tiny},
            xlabel={}, ylabel={},
        }
    }

    \begin{multicols}{4}
        \begin{plot}[zuordnung, title={\textbf{A}}]
            \addplot[domain=-2.5:2.5, red, thick, samples=100]{1/3*x^3-x};
        \end{plot}

        \begin{plot}[zuordnung, title={\textbf{B}}]
            \addplot[domain=-2.5:2.5, red, thick, samples=100]{-1/2*x^2+1};
        \end{plot}

        \begin{plot}[zuordnung, title={\textbf{C}}]
            \addplot[domain=-2.5:2.5, red, thick]{1/2*x-1};
        \end{plot}

        \begin{plot}[zuordnung, title={\textbf{D}}]
            \addplot[domain=-2.5:2.5, red, thick, samples=100]{-1/3*x^3+x};
        \end{plot}
    \end{multicols}

    \begin{multicols}{4}
        \begin{plot}[zuordnung, title={\textbf{1}}]
            \addplot[domain=-2.5:2.5, blue, thick, samples=100]{-x^2+1};
        \end{plot}

        \begin{plot}[zuordnung, title={\textbf{2}}]
            \addplot[domain=-2.5:2.5, blue, thick]{1/2};
        \end{plot}

        \begin{plot}[zuordnung, title={\textbf{3}}]
            \addplot[domain=-2.5:2.5, blue, thick, samples=100]{x^2-1};
        \end{plot}

        \begin{plot}[zuordnung, title={\textbf{4}}]
            \addplot[domain=-2.5:2.5, blue, thick]{-x};
        \end{plot}
    \end{multicols}

    \begin{center}
        \begin{tabularx}{.6\textwidth}{|c|X|X|X|X|}
            \hline
            Funktion & \centering A & \centering B & \centering C & \centering\arraybackslash D\\
            \hline
            Ableitungsfunktion & & & & \\
            \hline
        \end{tabularx}
    \end{center}

    \checkered[5cm]
\end{document}
